\section*{Capítulo \ref{ch:g9:metals-acids-corrosion} --- Los metales: reacciones con los ácidos y corrosión}

\begin{solution}{exo:g9:metals-acids-corrosion:1}
Hidrógeno: una astilla encendida en la boca del tubo da un pop agudo.
\end{solution}

\begin{solution}{exo:g9:metals-acids-corrosion:2}
Reaccionan el hierro, el cinc y el aluminio; el cobre y el oro, no.
\end{solution}

\begin{solution}{exo:g9:metals-acids-corrosion:3}
Oxígeno y agua, los dos a la vez.
\end{solution}

\begin{solution}{exo:g9:metals-acids-corrosion:4}
Un \omterm{def:g9:ions:ion}{ion} presente durante una reacción que no interviene en ella. Con el
ácido clorhídrico: el \omterm{def:g9:ions:ion}{ion} cloruro \ce{Cl-}.
\end{solution}

\begin{solution}{exo:g9:metals-acids-corrosion:5}
\ce{Mg(s) + 2H+(aq) -> Mg^{2+}(aq) + H2(g)}; carga $+2$ en cada miembro.
\end{solution}

\begin{solution}{exo:g9:metals-acids-corrosion:6}
Añadir hidróxido de sodio: un \omterm{def:g9:ions:precipitate}{precipitado} verde de \ce{Fe(OH)2} indica
\omterm{def:g9:ions:ion}{iones} hierro(II).
\end{solution}

\begin{solution}{exo:g9:metals-acids-corrosion:7}
Al hervirla se elimina el \omterm{def:g4:air-a-mixture-of-gases:air}{aire} disuelto en el agua, y el aceite impide
que el \omterm{def:g4:air-a-mixture-of-gases:air}{aire} vuelva a \omterm{def:g2:mixing-and-dissolving:dissolve}{disolverse}: el clavo tiene agua, pero no oxígeno.
\end{solution}

\begin{solution}{exo:g9:metals-acids-corrosion:8}
La sal disuelta en el agua que salpica el coche acelera la oxidación.
\end{solution}

\begin{solution}{exo:g9:metals-acids-corrosion:9}
La capa de óxido que se forma se adhiere al metal y lo aísla del \omterm{def:g4:air-a-mixture-of-gases:air}{aire}: el
ataque se detiene en la superficie.
\end{solution}

\begin{solution}{exo:g9:metals-acids-corrosion:10}
Pintarlo, aceitarlo o engrasarlo, galvanizarlo (o fabricarlo de acero
inoxidable).
\end{solution}

\begin{solution}{exo:g9:metals-acids-corrosion:11}
\ce{2Al(s) + 6H+(aq) -> 2Al^{3+}(aq) + 3H2(g)}: 2 Al y 6 H en cada
miembro, carga $+6$ en cada miembro. Para 200 \omterm{def:g7:atoms-and-molecules:atom}{átomos} de Al:
$200 \times \frac{3}{2} = 300$ \omterm{def:g7:atoms-and-molecules:molecule}{moléculas} de hidrógeno.
\end{solution}

\begin{solution}{exo:g9:metals-acids-corrosion:12}
El cinc de alrededor de la raya se corroe en lugar del acero: mientras
quede cinc cerca, el acero se libra.
\end{solution}

\begin{solution}{pb:g9:metals-acids-corrosion:1}
\textbf{1.} \ce{Zn(s) + 2H+(aq) -> Zn^{2+}(aq) + H2(g)}.

\textbf{2.} Hidrógeno: un pop agudo con una astilla encendida.

\textbf{3.} El hidróxido de sodio da un \omterm{def:g9:ions:precipitate}{precipitado} blanco de
\ce{Zn(OH)2}.

\textbf{4.} Los \omterm{def:g9:ions:ion}{iones} cloruro \ce{Cl-}.

\textbf{5.} $125.6 - 113.5 = \qty{12.1}{g}$.

\textbf{6.} Se desprende hidrógeno mientras el cinc reacciona con el
ácido; cuando ya no queda cinc, el burbujeo se detiene.

\textbf{7.} El hierro del acero también reaccionaría, dando hidrógeno e
\omterm{def:g9:ions:ion}{iones} \ce{Fe^{2+}}; el hidróxido de sodio daría entonces un \omterm{def:g9:ions:precipitate}{precipitado}
verde.

\textbf{8.} El cinc es atacado en lugar del acero y mantiene el agua y
el \omterm{def:g4:air-a-mixture-of-gases:air}{aire} lejos de él.

\textbf{9.} $2 \times 10.0 \times 10.0 = \qty{200}{cm^2}$.

\textbf{10.} $12.1 \div 7.13 \approx \qty{1.70}{cm^3}$.

\textbf{11.} $1.70 \div 200 = \qty{0.0085}{cm}$.

\textbf{12.} $0.0085 \times 10\,000 = \qty{85}{\micro m}$: la capa de
cinc tenía unos \qty{85}{\micro m} de espesor.
\end{solution}
